\section{Результаты и выводы}

\subsection{Задание 1.5. Сравнение результатов}

Записи речевого сигнала и характерные особенности осциллограмм фонем
носят индивидуальный характер и зависят от голосовых данных диктора
(значение $F_0$), особенностей произношения (длительность фонем,
сила артикуляции) и условий записи (уровень шума помещения,
характеристики микрофона).

Сравнение с результатами, полученными другими дикторами, показывает:
\begin{itemize}
	\item абсолютные значения $F_0$ изолированных гласных различаются
		у разных дикторов (для записанных фонем получено 119--205 Гц),
		однако \textbf{гармоническая структура спектра} и
		\textbf{квазипериодический характер} осциллограммы гласных
		сохраняются у всех дикторов;
	\item шипящие согласные у всех дикторов остаются
		шумоподобными (без гармоник на спектрограмме), при этом
		энергетические профили (распределение энергии по частотам)
		различаются в зависимости от артикуляции и голосовых
		особенностей;
	\item эффект коартикуляции в слогах «ба», «да», «ма», «на»
		наблюдается у всех дикторов --- это устойчивое свойство
		речевого сигнала, а не индивидуальная особенность;
	\item форма огибающих гласных и соотношение длительностей фонем
		варьируют от диктора к диктору, что делает осциллограммы
		конкретной записи неповторимыми при сохранении общих
		закономерностей.
\end{itemize}

\subsection{Задание 1.6. Общие выводы по работе}

В ходе выполнения работы были изучены фонетические закономерности
речевого сигнала. Основные результаты:
\begin{enumerate}
	\item Освоена запись речевых сигналов и регулировка уровней записи
		и воспроизведения; показано, что правильная установка уровня
		(пик около $-12$ дБFS, отсутствие клиппинга) обеспечивает
		качественный исходный материал для дальнейшего анализа.
	\item Исследовано влияние частоты дискретизации и разрядности:
		частота дискретизации определяет полосу частот (4 кГц ---
		разборчивость заметно ухудшается, 8 кГц --- телефонное
		качество, 16 кГц --- хорошее качество), разрядность 8 бит
		вносит слышимый шум квантования, но сохраняет разборчивость.
	\item Установлено, что гласные фонемы --- квазипериодические
		сигналы: они разборчиво звучат отдельно, а один
		повторяющийся период в циклическом режиме воспроизведения
		звучит как подобие исходной гласной. Буквы «я», «ё», «ю», «е»
		образуются парами фонем [й] + гласная.
	\item Показано, что шипящие согласные --- шумоподобные сигналы:
		они выделяются и прослушиваются отдельно, но отдельный период
		шипящей в циклическом режиме не звучит как эта фонема.
	\item Выявлен эффект коартикуляции: в слогах «фа», «ба», «да»,
		«ма» и т.\,д. согласные звучат только в контексте гласных ---
		их артикуляция накладывается на артикуляцию соседнего звука,
		и вырезать «чистую» согласную из слога невозможно.
\end{enumerate}

Таким образом, фонетическая структура речевого сигнала определяется
классом фонемы: гласные характеризуются периодической структурой и
формантным спектром, шипящие --- шумоподобным спектром в
характерной полосе, а в контексте слогов звучание согласных
существенно модифицируется за счёт коартикуляции. Полученные
закономерности являются основой для дальнейшего исследования
просодических характеристик (часть~2) и компиляционных методов
синтеза речи (часть~3).

Исходные звукозаписи (\texttt{sample1}---\texttt{sample5},
варианты частот дискретизации и разрядностей), скрипты анализа и
полученные рисунки сохранены в структуре каталогов лабораторной
работы и могут быть использованы при выполнении последующих частей.

\begin{thebibliography}{1}
	\bibitem{lobanov}
		Лобанов, Б.~М. Речевой интерфейс интеллектуальных систем :
		учеб. пособие / Б.~М.~Лобанов, О.~Е.~Елисеева ; под науч. ред.
		В.~В.~Голенкова. --- Минск : БГУИР, 2006. --- 152 с.
	\bibitem{rabiner}
		Рабинер, Л.~Р. Цифровая обработка речевых сигналов /
		Л.~Р.~Рабинер, Р.~В.~Шафер ; пер. с англ. --- М. : Радио и
		связь, 1981. --- 495 с.
\end{thebibliography}
